

\subsection{Motivating Example}
                                                  
Below we show a toy example problem that illustrates
the need for MEGA. We will use this simple example throughout
this paper to explain various aspects of our approach
and to contrast our approach with others. 
This example models an admission concept that
consists of a small number of boolean predicates.
In Section~\ref{sec-expimage} we
show that MEGA works well to learn complex query concepts 
of fuzzy predicates
for a high-dimensional image dataset.

Suppose Jane plans to apply to
a graduate school. Before filling out the
forms and paying the application fees, she would like
to estimate her chances of being admitted. Since
she does not know the admission criteria, 
she decides to learn the admission concept by 
induction\footnote{
At first sight, this example may seem very different from, say,
an image search scenario, where a user queries similar
images by example(s). But if we treat the admission officer
as the user who knows what he/she likes and who can, accordingly, 
label an item as {\it true} or {\it false}, and if
we treat Jane as the search engine who tries to
find out what the admission officer thinks, then
it is evident that this example represents a typical
search scenario.
}.
She calls up a few friends who applied
last year and obtains
the information shown in Table~\ref{tab-ex1}.

\begin{table}[ht]
\begin{footnotesize}
\begin{center}
\begin{tabular}{|c|c|c|c|c|} \hline
{\it Name} & {\it GPA} & {\it GRE} & {\it Has Publications?} & {\it Was Admitted?} \\ \hline \hline
{\it Joe}   & {\it high} & {\it high} & {\it false} & {\it true}  \\
{\it Mary}  & {\it high} & {\it low}  & {\it true}  & {\it true}  \\
{\it Emily} & {\it high} & {\it low}  & {\it true}  & {\it true}  \\
{\it Lulu}  & {\it high} & {\it high} & {\it true}  & {\it true}  \\
{\it Anna}  & {\it low}  & {\it low}  & {\it true}  & {\it false} \\
{\it Peter} & {\it low}  & {\it high} & {\it false} & {\it false} \\
{\it Mike}  & {\it high} & {\it low}  & {\it false} & {\it false} \\
{\it Pica}  & {\it low}  & {\it low}  & {\it false} & {\it false} \\
\hline
\hline
\end{tabular} 
\end{center}
\end{footnotesize}
\caption{Admission Samples.}
\label{tab-ex1}
\end{table}

If we look at the GRE scores in the table, we see that students with 
either high or low GRE scores were admitted, and both kinds were
rejected. Hence, we may 
conclude that the GRE is irrelevant in the admission process. 
Likewise, one's publication record does not affect admission
acceptance, nor does having a high GPA.  
It may appear that the admission decision is entirely
random. However, the 
graduate school actually uses a combination of
reasonable criteria: it requires a
high GPA and either a high GRE or publications. In other words,
$Admission~:~ GPA = high \wedge (GRE = high \vee Publications = true)$.

An obvious question arises: ``Are all the samples in Table~\ref{tab-ex1}
equally useful for learning the target concept?''
% and, ``What features in the table are relevant to the learning task?''
%\begin{jlist}
%\item Are all samples equally useful?
Apparently not, for several
reasons.  First, it seems that $Pica$'s record may not be useful
since she was unlikely to be admitted (i.e., her record is unlikely
to be labeled positive). Second, both Emily and Mary have the 
same record, so one of these two records can be redundant. Third, Lulu's 
record is perfect and hence does not provide additional insight for
learning the admission criteria.  This example indicates that choosing 
samples randomly may not produce useful information for learning a 
target concept. 
In this paper, we will show how and why MEGA can quickly learn
a target concept like the example of
admission criteria, whereas other methods may not.
%\item Are all features relevant? To determine relevancy, we
%examine the features in the table.
%The feature {\it ``Gender''} does not seem to
%be relevant to graduate admissions.
%The presence of irrelevant features
%can slow down concept learning exponentially \cite{LI93,LS97}.
%\end{jlist}




\section{The MEGA Algorithm}
\label{sec-learn}

%In this section, we describe
%how our learning algorithm, MEGA, works.
%In Section~\ref{sec-db-issues}, we discuss how MEGA deals
%with very large database issues such as high dimensional data and
%very large datasets.

Our learning algorithm consists of two major parts:
\begin{jlist}
%\item Initialization: Provide users with a reasonable way to
%convey initial hints to the system.
\item Refinement: Refine the query concept based on 
positive-labeled instances. The refinement step 
is carefully designed to tolerate noisy data.
\item Sampling: Refine the sampling space based
on negative-labeled instances and select samples
judiciously for expediting the learning process.
\end{jlist}

\subsection{Refinement}

We use Valiant's learning algorithm \cite{V84}
as the starting point to refine a $k$-CNF concept. 
We extend the algorithm to
\begin{nlist}
\item {\bf 1}: Handle the 
fuzzy membership functions (described in Section~\ref{sec-formulate}), 
\item {\bf 2}: Select samples judiciously to
expedite the learning process (Section~\ref{sec-sampling}), and
\item {\bf 3}: Tolerate user errors (Section~\ref{sec-vote}).
\end{nlist}

\begin{figure}[ht]
\epsfysize=100pt
\centerline{\epsffile{fig-qccs.eps}}
\vspace{-.1in}
\caption{MEGA's Sampling Space: $CCS \wedge \overline{QCS}$.}
\vspace{-.1in}
\label{fig-qccs}
\end{figure}
               
MEGA initializes
the query concept space ($QCS$) as the most specific $k$-CNF
and the candidate concept space ($CCS$) as the most general $k$-DNF.
%The $QCS$ starts as the most specific concept
%and the $CCS$ as the most general concept.
The target concept
that MEGA learns is more general than the initial
$QCS$ and more specific than the initial $CCS$. 
It will become evident shortly that MEGA learns the $QCS$,
while at the same time refining the $CCS$ to delimit
the boundary of the sampling space. (The
shaded area in Figure~\ref{fig-qccs} shows the sampling space
between the $QCS$ and the $CCS$).

Figure~\ref{fig-learner} depicts the learning algorithm.
The first step initializes $QCS$ and $CCS$.
The second step describes how a query concept is
refined through relevance feedback. This refinement
step consists of the four substeps shown.

\begin{figure}[htbp]
%\begin{footnotesize}
{\bf Algorithm MEGA}
\begin{nlist} 
\item {\it Input}: $U$, $K_c$, $K_d$, $K_{\alpha}$;
\item {\it Output}: $QCS$;
\item {\it Procedure calls}: $f()$, $Vote()$, $Sample()$, $Feedback()$, 
                            $Collapsed?()$, $Converged?()$;
\item {\it Variables}: $c$, $d$, $u$, $y$, $P_{t_i|y_j}$, $P_{t_i|y}$;
\item {\bf Begin}
\item {\bf 1} Initialize the version space
 \begin{nlist}
 \item $QCS \leftarrow \{d_1, d_2, \dots\}$; $CCS \leftarrow \{c_1, c_2, \dots\}$;
 \end{nlist}
\item {\bf 2} Refine query concept via relevance feedback
  \begin{nlist}
   \item {While} ($not~Collapsed?(QCS, CCS)~and~not~Converged?(QCS, CCS)$) 
   \item {\bf 2.a} Select informative samples
   \begin{nlist}
   \item $u \leftarrow Sample(QCS, CCS, U, K_{\alpha})$; 
   \end{nlist}
   \item {\bf 2.b} Solicit user feedback
   \begin{nlist}
   \item For each $u_i \in u$
    \begin{nlist}
    \item $y_i \leftarrow Feedback(u_i)$; 
    \end{nlist}
   \end{nlist}
   \item {\bf 2.c} Reduce irrelevant terms
   \begin{nlist}
   \item {\bf 2.c.1} Refine $k$-CNF using positive samples

   \begin{nlist}
   \item For each $d_i \in QCS$
    \begin{nlist}
     \item For each $y_j \in y^+$
     \begin{nlist}
     \item $P_{d_i|y_j} \leftarrow f(d_i, O_j, QCS)$;
     \end{nlist}
    \item $P_{d_i|y^+} \leftarrow Vote(y^+, P_{d_i | y_j \in y^+}, \gamma)$;
    \item If ($P_{d_i|y^+} > K_d$)
     \begin{nlist}
     \item $QCS \leftarrow QCS - \{d_j\}$;
     \end{nlist}
    \end{nlist}
   \end{nlist}

   \item {\bf 2.c.2} Refine $k$-DNF using negative samples
   \begin{nlist}
   \item For each $c_i \in CCS$
    \begin{nlist}
     \item For each $y_j \in y^-$
     \begin{nlist}
     \item $P_{c_i|y_j} \leftarrow f(c_i, O_j, CCS)$;
     \end{nlist}
    \item $P_{c_i|y^-} \leftarrow Vote(y^-, P_{c_i | y_j \in y^-}, \gamma)$;
    \item If ($P_{c_i|y^-} > K_c$)
     \begin{nlist}
     \item $CCS \leftarrow CCS - \{c_j\}$;
     \end{nlist}
    \end{nlist}
   \end{nlist}
  \end{nlist}

  \item {\bf 2.d} Bookkeeping
   \begin{nlist}
   \item $U \leftarrow U - u$;
   \end{nlist}
\end{nlist}

\item {\bf 3} Return query concept
\begin{nlist}
  \item Output $QCS$; 
\end{nlist}
\item {\bf End}
\end{nlist}  
\caption{Algorithm MEGA.}
\label{fig-learner}
%\end{footnotesize}
\end{figure}

\begin{nlist}
\item {\bf Step 2.a}: 
MEGA Generates
sample set $u$ from the unlabeled pool $U$.
This step passes $QCS$, $CCS$ and $U$ to procedure {\it Sample}
to generate $K_{\alpha}$ samples.
We will discuss procedure {\it Sample} in Section~\ref{sec-sampling}.

\item {\bf Step 2.b}: MEGA solicits user feedback. A user marks
an object positive if the object fits his/her query
concept.  An unmarked object is considered negative. 

\item {\bf Step 2.c}: 
This is MEGA's term reduction step. It
refines both $QCS$ and $CCS$. 

{Step 2.c.1} refines $QCS$.  For each disjunctive term in the $QCS$, 
MEGA examines each positive-labeled instance
and uses function $f$ to compute the probability that
the disjunctive term should be eliminated.
MEGA then calls procedure {\it Vote} to tally the votes among the
positive-labeled instances and compares the vote with threshold $K_d$
to decide whether that disjunctive term is to be removed. 
If sufficient numbers of positive-labeled instances contradict the
$QCS$ with respect to a disjunctive term (i.e., if the
threshold is exceeded), the term is removed from the $QCS$.
We discuss procedure {\it Vote}, which decides how
aggressive MEGA is in eliminating terms, in Section~\ref{sec-vote}.

{Step 2.c.2} refines $CCS$.
Similar to the steps in 2.c.1, for each conjunctive term in
the $CCS$, MEGA examines each negative-labeled instance,
and uses function $f$ to compute the probability that
the conjunctive term should be eliminated.
MEGA then calls procedure {\it Vote}
to tally the votes among the
negative-labeled instances.
Then it compares the vote with threshold $K_c$
to decide whether that conjunctive term is to be removed.
If sufficient numbers of negative-labeled 
instances satisfy the $CCS$ with respect to
a conjunctive term, the term is removed from the $CCS$.

\item {\bf Step 2.d}: MEGA performs bookkeeping by reducing the
unlabeled pool.
\end{nlist}

The refinement step terminates when the learning
process converges to the target concept ($Converged?=true$) or the
concept is collapsed ($Collapsed?=true$).  ($Converged?$ and
$Collapsed?$ are defined below.)
In practice, the refinement stops when no unlabeled instance
$u$ can be found between the $QCS$ and the $CCS$.

\vspace{.1in}
\sspace
\noindent
{\bf Definition 3:} $Converged?(QCS, CCS)$
\[
Converged?(QCS, CCS) \leftarrow~true~if~CCS == QCS;~~false~otherwise.
\]

\noindent
{\bf Definition 4:} $Collapsed?(QCS, CCS)$
\[
Collapsed?(QCS, CCS) \leftarrow~true~if~CCS \in QCS;~~false~otherwise.
\]
\nspace

\subsection*{Remarks}

Notice that MEGA does not require seeding a
query with relevant images.  MEGA intelligently generates
images to seek for user feedback.
If all of the images
generated in the first round of its relevance feedback are irrelevant,
it will use these irrelevant images to reduce
the sampling space
%the set of potentially relevant images substantially
so that a relevant image will be
generated in the next round with high probability.
                                                      

\subsection{Sampling}
\label{sec-sampling}

MEGA invokes procedure {\it Sample}
to select the next $K_{\alpha}$ unlabeled instances
%\footnote{
%In this study we focus our discussion on sample selection
%criteria rather than on access methods for finding samples
%on disk. We are currently enhancing a high-dimensional
%indexing scheme that we have developed \cite{Lch00} 
%to achieve efficient sample retrieval.} 
to ask for user feedback. From the 
college-admission example presented in Section~\ref{sec-intro}, we
learn that if we would like to reduce our work (i.e.,
call a minimum number of friends), we should choose
our samples judiciously. 
But, what constitutes a good sample?  We know that we
learn nothing from a sample if
\begin{jlist}
\item It agrees with the concept in all terms.
\item It has the same features as another sample.
\item It is unlikely to be labeled positive.
\end{jlist}

To make sure that a sample is useful, MEGA employs two strategies:
\begin{nlist}
\item {\bf 1}. Bounding the sample space:
MEGA avoids choosing useless unlabeled instances by using
the $CCS$ and $QCS$ to delimit the sampling boundary. 
%Figure~\ref{fig-qccs} shows the candidate space where MEGA
%selects unlabeled instances is between $CCS$ and $QCS$
%(the shaded area in the figure).
\item {\bf 2}. Maximizing the usefulness of a sample:
MEGA chooses a sample that
shall remove the maximum expected number of disjunctive
terms in QCS. In other words, MEGA chooses a sample
that can maximize the expected generalization of the concept.
\end{nlist}

It may appear that
if we pick a sample that has more dissimilar
disjunctions (compared to the $QCS$), we may have
a better chance of eliminating more disjunctive terms. 
This is, however, not true.
A sample must be labeled by the user as positive to be
useful for refining $QCS$. Unfortunately, a sample with
more disjunctions that are dissimilar
to the target concept is less likely to be
labeled positive.  Therefore, in choosing a sample,
there is a tradeoff between those with more contradictory
terms and those more likely to be labeled positive.

Let $\Psi$ denote the number of disjunctions remaining in the concept.
The number of disjunctions that can be eliminated in the current round 
of sampling (denoted as $\psi$) is between zero and $\Psi$.
For a sample contradicting exactly $\psi$ disjunctive terms in QCS,
we can
write the probability of its being labeled positive as $P_e(\psi)$.
$P_e(\psi)$ is a monotonically decreasing function of $\psi$.

Our MEGA scheme attempts to find the $\psi$ that
can eliminate the maximum expected number
of disjunctive terms, given a sample.
MEGA's objective function can be written as
\begin{equation}
\psi^{*} = argmax_{\psi} E(\psi) = argmax_{\psi} (\psi \times P_e(\psi)).
\label{eq-mega}
\end{equation}
To solve $\psi^{*}$, we must know $P_e(\psi)$, which
can be estimated by the two methods below: {\it probabilistic estimation}
and {\it empirical estimation}.

\subsection{Probabilistic Estimation}

We first consider how to estimate $\psi^*$ using a
probability model.
As we have seen in the college-admission analogy, if a sample contradicts
more disjunctive terms, it is more likely to be labeled negative
(i.e., less likely to be labeled positive).
For example, a sample that contradicts term $T_1$ is
labeled negative only if $T_1$ is in the user's query concept.
A sample that contradicts both terms
$T_1$ and $T_2$ is labeled negative if
either $T_1$ or $T_2$ is in the user's query concept.

Formally, let random variable $\Phi_i$ be 1 if $T_i$ is in
the concept and $0$ otherwise.  For simplicity, let us assume 
that the $\Phi_i$'s are iid (independent and identically distributed),
and the probability of $\Phi_i$ being $1$ is $p$ ($0<p<1$).
For a sample contradicting $\psi$ disjunctive terms, it
is marked positive only when
none of these $\psi$ terms appears in the user's query concept. 
This probability is $(1-p)^\psi$.  
If we substitute $P_e(\psi)$ by $(1-p)^\psi$
on the right-hand side of Equation~\ref{eq-mega}, we get
\[
E(\psi) = \psi(1 - p)^\psi.
\]
If we take the derivative of 
$E(\psi)$, we can find the optimal
$\psi$ value, denoted by $\psi^{*}$:
\[
\psi^{*} = \Psi,~if~\frac{1}{ln \frac{1}{1-p}} > \Psi,~~~
\psi^{*} = \frac{1}{ln \frac{1}{1-p}},~ otherwise.
\]

%We can round the value up to the nearest integer, or
%
%%\noindent
%$\psi^{*} = \Psi$, if $\lceil \frac{1}{ln \frac{1}{1-p}} \rceil > \Psi$,
%
%%\noindent
%$\psi^{*} = \lceil \frac{1}{ln \frac{1}{1-p}} \rceil$, otherwise.
%Your sampling method assumes that each clause in the current concept is
%equally likely to be eliminated - clearly a very big assumption. From
%what I understand, the sampling method only gives you the number of
%clauses that it wants to eliminate rather than the specific clauses it 
%thinks it can eliminate. I wonder if there is a way to make the sampling
%algorithm take into account the structure of specific clauses in the
%current concept.

Of course, it may be too strong an assumption that the 
probability $p$ of all disjunctions is iid.
However, we do not need a precise estimation here for 
the following two reasons:
\begin{enumerate}
\item Precise estimation may not be feasible
and can be computationally intensive.
\item An approximate estimation is sufficient for bootstrapping.
Once the system is up and running for a while and collects
enough data, it can empirically
estimate $P_e(\psi)$ using its past experience.  We 
discuss this process next.
\end{enumerate}

\subsection{Empirical estimation}

MEGA can estimate
$P_e(\psi)$ based on its past experience. 
For each sample it presents, it records how many disjunctions
the sample contradicts with respect to the query concept and whether
the sample is labeled positive.  Once the system collects
a sufficient amount of data, we can estimate $P_e(\psi)$
empirically. We then pick the $\psi^{*}$ that can
eliminate the maximum expected number of disjunctive terms. 

Again, a reasonable approach to estimate $P_e(\psi)$ is to use
probabilistic estimation when the system first
starts and then switches to empirical estimation when the
system has collected sufficient data.
In such a case there can
be an abrupt transition when MEGA switches
from one estimation to the other, since the two
estimates of $P_e(\psi)$ may differ substantially. 
This may lead to a sudden change in behavior of the sampling 
component of the active learner.  To remedy this problem, 
we employ a Bayesian smoothing approach.
Essentially the probabilistic estimation is the prior guess at the
distribution over $\psi$ and the empirical approach is the guess based 
purely on the data that MEGA has seen so far. 
The Bayesian approach
combines both of these guesses in a principled manner. 
Before we start, we imagine that we have seen a number of
samples of $\psi$. After refinement iteration
we gather new samples for $\psi$; then we
add them to our current samples and adjust $P_e(\psi)$. 

For example, before we start, we assume that 
we have already seen samples with $\psi= 1$ being labeled positive
three out of five times and samples with
$\psi = 2$ being labeled positive seven out of $20$ times.
In other words, we
have successfully eliminated $\psi = 1$ term three times out of
five, and we have successfully eliminated $\psi = 2$ terms
$7$ times out of $20$.  Thus initially 
$P_e(\psi = 1) = 3/5 = 0.6$ and $P_e(\psi = 2) = 7/20 = 0.35$.
Now suppose we do a query in which we observe a
sample with $\psi = 2$ being labeled positive. Then our 
new distribution is
$P_e(\psi = 1) = 3/5$ and $P_e(\psi = 2) = 8/21$.
We continue in this manner.  At first the prior assumption 
has quite an effect on our guess about the
distribution. The more imaginary samples we have in our prior 
assumption, the larger its effect. 
For instance, if we assume that $\psi = 1$ being
labeled positive $30$ out of $50$ times and that
$\psi = 2$ being labeled positive $70$ out of $200$ times,
it takes more real samples to change $P_e(\psi)$.
With time, the more real samples we get, the less the
effect of the prior assumption becomes, 
until eventually it has virtually no
effect, and the observed data dominate the expression. 
This procedure gives us a smooth transition between the ``probabilistic'' 
and the ``empirical'' methods.

%(Formally, in this example we have placed a Beta(3,2) prior on the
%parameter $\theta_a$ where $\theta_a$ corresponds to P(A=true). We have also
%placed a Beta(7,20) prior on $\theta_b$. We then have done regular Bayesian  
%updating.
%
%For a more complete explanation take a look at the first couple of
%sections of David Heckerman's tutorial on learning with Bayesian
%Networks:
%http://www.research.microsoft.com/scripts/pubdb/pubsasp.asp?recordID=63 )
                     
%\begin{figure}
%\begin{center}
%\begin{tabular}{lr}
%\epsfxsize=120pt
%\epsfysize=80pt
%\subfigure[Random]{\epsfbox{fig-mega-random.eps}} 
%\hspace{.2in}
%\epsfxsize=120pt
%\epsfysize=80pt
% \subfigure[Guarded Random]{\epsfbox{fig-mega-G-random.eps}} \\     
%\epsfxsize=120pt
%\epsfysize=80pt
% \subfigure[MEGA]{\epsfbox{fig-mega-mega.eps}} 
%\hspace{.2in}
%\epsfxsize=120pt
%\epsfysize=80pt
% \subfigure[Aggressive]{\epsfbox{fig-mega-aggresive.eps}} \\
%\end{tabular}
%\end{center}
%\vspace{-.2in}
%\caption{Sampling Schemes.}
%\vspace{-.1in}
%\label{fig-mega}
%\end{figure}      
%

\subsection*{Remarks}

Once the criteria for selecting samples is determined, we
use Clindex, a clustering/indexing combined indexer to retrieve
the most informative samples.  We discuss this high-dimensional
access method in detail in \cite{Lch00}.  

\subsection*{Sampling Example}

Let us revisit the college-admission analogy to explain
how MEGA works. Suppose Joe is interested in applying for 
a graduate school.
He knows the profiles of a few students in his 
senior class, but he does not know if they were admitted
to that particular school. 
According to MEGA, we write the initial $CCS$ and $QCS$ in
$2$-CNF and $2$-DNF, respectively, as

\sspace
\[
QCS = (GRE = high)~ \wedge ~(GPA = high)~ \wedge~ (Publications = true) 
~\wedge~ (GRE = high \vee GPA = high) 
\]
\[
~\wedge~ (Publications = true \vee GPA = high) ~\wedge~
(GRE = high \vee Publications = true).
\]
\vspace{-.1in}
\[
CCS = (GRE = high)~ \vee ~(GPA = high)~ \vee~ (Publications = true) 
~\vee~ (GRE = high \wedge GPA = high) 
\]
\[
~\vee~ (Publications = true \wedge GPA = high) ~\vee~ 
(GRE = high \wedge Publications = true).
\]
\nspace

Suppose $\psi^{*}$ is one. 
Joe starts by calling his friends
whose ``profile'' fails by exactly one disjunctive term.
Joe calls three people and two give him a positive answer,
as shown in Table~\ref{tab-round1}.

\begin{table}[ht]
\begin{footnotesize}
\begin{center}
\begin{tabular}{|c|c|c|c|c|c|} \hline
{\it Round $\#$} & {\it Name} & {\it GPA} & {\it GRE} & {\it Has Publications?} & {\it Was Admitted?}\\ \hline
\hline
{\it 1st} & {\it Joe} & {\it high} & {\it high} & {\it false} & {\it true} \\
{} & {\it Emily} & {\it high} & {\it low} & {\it true} & {\it true} \\
{} & {\it Dora} & {\it low} & {\it high} & {\it true} & {\it false} \\
\hline \hline
{\it 2nd} & {\it Kevin} & {\it high} & {\it low} & {\it false} & {\it false} \\
\hline \hline
\end{tabular}
\end{center}
\end{footnotesize}
\caption{MEGA Sampling Rounds.}
\label{tab-round1}
\end{table}
                             
Based on the feedback, Joe can generalize the concept to
$QCS = (GPA = high)~\wedge~ (Publications = true \vee GPA = high) ~\wedge~
(GRE = high \vee Publications = true)~\wedge~
(GPA = high \vee GRE = high)$.
At the same time, the $CCS$ is shrunk to
$CCS = (GPA = high)~ \vee~
(GRE = high \wedge GPA = high)
~\vee~ (Publications = true \wedge GPA = high)$.

In the second round, Joe attempts to call friends to
see if any of the remaining terms can be removed. He calls
Kevin, whose profile is listed in the table.
Since this sample is labeled negative, the $QCS$ is not changed.
But the $CCS$ is reduced to 
$(GRE = high \wedge GPA = high)
~\vee~ (Publications = true \wedge GPA = high)$.

Simplifying and rewriting
both $QCS$ and $CCS$ gives us the following identical
expression:
\[
QCS = (GPA = high) ~\wedge~ (GRE = high \vee Publications = true).
\]
The concept converges and the
refinement terminates at this point. We have learned the
admission criterion --- a high GPA and either a high GRE or
publications $\Box$

\subsection{Procedure Vote}
\label{sec-vote}

In the previous example, we assume that all samples are noise-free.
This assumption may not be realistic. There can be two sources of
noise:

\begin{jlist}
\item Model bias: The target concept may not be in $k$-CNF.
\item User errors: A user may label some positive instances negative
and vice versa.
\end{jlist}

\begin{figure}[thbp]
{\bf Procedure Vote}
\begin{nlist}
\item {\it Input}: $y$, $P_{t_i| y_j \in y}$, $\gamma$;
\item {\it Output}: $P_{t_i|y}$;
%\item {\it Procedure call}: $\omega$(); /* weighting function */
\item {\bf Begin}
\begin{nlist}

\item Sort $P_{t_i|y_j}$ in the descending order;
\item Return the $\gamma^{th}$ highest $P_{t_i|y_j}$;
\end{nlist}
\item {\bf End}
\end{nlist}
\caption{Procedure Vote.}
\end{figure}        

Procedure {\it Vote} controls the strictness of voting using
$\gamma$.  The larger the value of $\gamma$ is, the more strict the
voting is and therefore the harder it is to eliminate a term.
When the noise level is high, we have less confidence in the correctness
of user feedback. Thus, we want to be more cautious 
about eliminating a term.
Being more cautious means increasing $\gamma$.
Increasing $\gamma$, however, makes the learning process
converge more slowly.  To learn a concept when noise is present,
one has to buy accuracy with time.  
In Section~\ref{sec-errors} we show how the voting threshold $\gamma$ affects
search accuracy and convergence speed through our experimental results,
and we suggest heuristics for selecting a proper $\gamma$.

